% Einheit 5 – Potenzfunktionen mit natürlichem Exponenten
\setcounter{aufgabe}{42}
\hypertarget{e5}{}\einheitenkopf*{Einheit 5 von 5 · Potenzfunktionen mit natürlichem Exponenten}
\zweigzeile{Hier lernst du Funktionen wie $y = x^3$ und $y = x^4$ kennen: Wertetabelle, Graph und Eigenschaften \,·\, neu in diesem Jahr (am Gymnasium schon seit Klasse 9) \,·\, keine P10-Aufgabe \,·\, baut auf: Einheit 3, Normalparabel, Punkte im Koordinatensystem}

\begin{aufgabe}{Ich kann Wertetabellen zu Potenzfunktionen ausfüllen. Berechne die fehlenden Werte.}
\beispiel{y = x^3,\ x = -2: \quad y &= (-2)^3 = (-2) \cdot (-2) \cdot (-2) = -8}
\begin{teile}
\teil $y = x^3$
\end{teile}
\wertetabelle{x}{y}{0,1,2,3}
\begin{teile}
\teil $y = x^3$, jetzt mit negativen Zahlen
\end{teile}
\wertetabelle{x}{y}{-3,-2,-1}
\begin{teile}
\teil $y = x^4$
\end{teile}
\wertetabelle{x}{y}{-2,-1,0,1,2}
\begin{teile}
\teil $y = x^4$, auch mit Dezimalzahlen
\end{teile}
\wertetabelle{x}{y}{-0{,}5,0{,}5,3}
\end{aufgabe}

\begin{aufgabe}{Ich kann Funktionswerte von Potenzfunktionen berechnen. Berechne.}
\beispiel{f(x) = x^4: \quad f(3) &= 3 \cdot 3 \cdot 3 \cdot 3 = 81}
\begin{teilezwei}
\tz $f(x) = x^3$: \ \feld{f(2)} & \tz $f(x) = x^4$: \ \feld{f(2)} \\
\tz $f(x) = x^3$: \ \feld{f(3)} & \tz $f(x) = x^4$: \ \feld{f(10)} \\
\tz $f(x) = x^3$: \ \feld{f(-5)} & \tz $f(x) = x^4$: \ \feld{f(-2)} \\
\tz $f(x) = 2x^3$: \ \feld{f(-2)} & \tz $f(x) = x^3$: \ \feld{f(0{,}5)} \\
\end{teilezwei}
Prüfe, ob der Punkt auf dem Graphen liegt.
\begin{teilezwei}
\tz $P(-2\,|\,16)$, \ $y = x^4$: \ \janein & \tz $Q(3\,|\,{-27})$, \ $y = x^3$: \ \janein \\
\end{teilezwei}
Bestimme alle $x$, für die die Gleichung gilt.
\begin{teilezwei}
\tz $x^3 = 216$: \ \leerfeld & \tz $x^4 = 81$: \ \leerfeld \\
\end{teilezwei}
\end{aufgabe}

\begin{aufgabe}{Ich kann den Graphen einer Potenzfunktion zeichnen. Nutze deine Wertetabellen aus Nr.~43 und zeichne die Graphen in das Koordinatensystem.}
\begin{teile}
\teil $y = x^2$ (Normalparabel) \quad \teil $y = x^3$ \quad \teil $y = x^4$
\end{teile}

\begin{ksys}[xmin=-2,xmax=2,ymin=-8,ymax=16,xstep=0.5,ystep=2]
\end{ksys}

\end{aufgabe}

\begin{aufgabe}{Ich kann an der Gleichung erkennen, wie der Graph einer Potenzfunktion verläuft. Kreuze an, wozu der Graph symmetrisch ist.}
\begin{teile}
\teil $y = x^4$ \hfill \kreuz{zur $y$-Achse}\kreuz{zum Ursprung}
\teil $y = x^5$ \hfill \kreuz{zur $y$-Achse}\kreuz{zum Ursprung}
\teil $y = x^6$ \hfill \kreuz{zur $y$-Achse}\kreuz{zum Ursprung}
\teil $y = x^3$ \hfill \kreuz{zur $y$-Achse}\kreuz{zum Ursprung}
\end{teile}
Vergleiche mit dem Graphen von $y = x^4$. Kreuze an.
\begin{teile}
\teil $y = 3x^4$ \hfill \kreuz{schmaler}\kreuz{breiter}
\teil $y = 0{,}25x^4$ \hfill \kreuz{schmaler}\kreuz{breiter}
\teil Beschreibe, wie der Graph von $y = -x^4$ aus dem Graphen von $y = x^4$ entsteht.
\end{teile}
\end{aufgabe}

\begin{aufgabe}{Ich kann Gleichungen und Graphen von Potenzfunktionen einander zuordnen. Gib an, welcher Graph zur Gleichung gehört.}

\begin{ksys}[xmin=-2,xmax=2,ymin=-4,ymax=4,ablesen]
\funktion[1.25]{\x^4}{\mathrm{I}}
\end{ksys}\hspace{0.8cm}\begin{ksys}[xmin=-2,xmax=2,ymin=-4,ymax=4,ablesen]
\funktion[1.2]{\x^3}{\mathrm{II}}
\end{ksys}\hspace{0.8cm}\begin{ksys}[xmin=-2,xmax=2,ymin=-4,ymax=4,ablesen]
\funktion[1]{-\x^2}{\mathrm{III}}
\end{ksys}\hspace{0.8cm}\begin{ksys}[xmin=-2,xmax=2,ymin=-4,ymax=4,ablesen]
\funktion[1.5]{0.5*\x^3}{\mathrm{IV}}
\end{ksys}

\begin{teilezwei}
\tz $y = x^3$: \ Graph \leerfeld & \tz $y = -x^2$: \ Graph \leerfeld \\
\tz $y = x^4$: \ Graph \leerfeld & \tz $y = 0{,}5x^3$: \ Graph \leerfeld \\
\end{teilezwei}
\begin{teile}
\teil Welcher der vier Graphen hat dieselbe Form wie der Graph von $y = x^6$? Kreuze an. \\ \kreuz{I}\kreuz{II}\kreuz{III}\kreuz{IV}
\end{teile}
\end{aufgabe}

\begin{aufgabe}{Ich kann unterscheiden, ob die Variable in der Basis oder im Exponenten steht. Kreuze an.}
\begin{teile}
\teil $y = x^3$ \hfill \kreuz{Potenzfunktion}\kreuz{Exponentialfunktion}
\teil $y = 3^x$ \hfill \kreuz{Potenzfunktion}\kreuz{Exponentialfunktion}
\teil $y = 2x^5$ \hfill \kreuz{Potenzfunktion}\kreuz{Exponentialfunktion}
\teil $N(t) = 50 \cdot 1{,}1^t$ \hfill \kreuz{Potenzfunktion}\kreuz{Exponentialfunktion}
\teil $y = 0{,}5^x$ \hfill \kreuz{Potenzfunktion}\kreuz{Exponentialfunktion}
\teil $y = x^{10}$ \hfill \kreuz{Potenzfunktion}\kreuz{Exponentialfunktion}
\teil Berechne $x^3$ und $3^x$ für $x = 2$ und für $x = 5$. Welcher Term wächst für große $x$ schneller?
\end{teile}
\end{aufgabe}

\begin{aufgabe}{Ich finde den Fehler beim Berechnen von Funktionswerten. Tom soll für $f(x) = x^3$ den Wert $f(-2)$ berechnen. Er rechnet so:}
\rechnung{f(-2) &= -2 \cdot 3 = -6}
\begin{teile}
\teil Finde den Fehler, benenne ihn und korrigiere die Rechnung.
\teil Berechne selbst für $f(x) = x^3$ den Wert $f(-3)$. \quad \leerfeld
\end{teile}
Sara soll für $f(x) = -x^4$ den Wert $f(2)$ berechnen. Sie rechnet so:
\rechnung{f(2) &= (-2)^4 = 16}
\begin{teile}
\teil Finde den Fehler, benenne ihn und korrigiere die Rechnung.
\teil Berechne selbst für $f(x) = -x^4$ den Wert $f(3)$. \quad \leerfeld
\end{teile}
\end{aufgabe}

\begin{aufgabe}{Ich kann Eigenschaften von Potenzfunktionen begründen. Entscheide ohne Rechnung und begründe.}
\begin{teile}
\teil Warum wird $y = x^4$ nie negativ?
\teil Warum verläuft der Graph von $y = x^3$ links vom Ursprung unterhalb der $x$-Achse?
\end{teile}
\end{aufgabe}

\begin{aufgabe}{Ich kann mit einer Potenzfunktion eine Sachaufgabe lösen. Ein Würfel hat die Kantenlänge $x$ in cm. Sein Volumen in $\text{cm}^3$ ist $V = x^3$.}
\begin{teile}
\teil Berechne das Volumen eines Würfels mit der Kantenlänge $7\,\text{cm}$. \quad \leerfeld[$\text{cm}^3$]
\teil Ein Würfel soll ein Volumen von $1\,331\,\text{cm}^3$ haben. Wie lang ist seine Kante? \quad \leerfeld[cm]
\teil Die Kante eines Würfels wird verdoppelt. Wie verändert sich sein Volumen? Begründe.
\end{teile}
\end{aufgabe}
